\chapter*{Acknowledgements}
\renewcommand{\baselinestretch}{\linespacing}
\small\normalsize

I would like to thank my supervisor Christoph Flamm for the opportunity and his patience.
For tolerating my occasional slow grasp of concepts and for all the support I received, I am grateful.

Gratitude is also extended to the entire TBI working group for their support and for creating a stimulating environment. Special thanks to Manuel Uhlir for getting me started and for the many helpful discussions and suggestions along the way.

Heartfelt thanks are due to my mother and grandmother, whose constant faith in me has carried me through many challenges. I also wish to acknowledge my grandfather for the memorable time I was able to share with him.

Finally, thank you, Aila, for your warmth and kindness at every step of this journey.

\newpage
\pdfbookmark[0]{Contents}{contents_bookmark}
\renewcommand{\baselinestretch}{1.3}\normalsize
\tableofcontents
\renewcommand{\baselinestretch}{1.5}\normalsize
\clearpage
\chapter*{Abstract}
%min. 500 Zeichen inkl. Leerzeichen

Mass spectrometry is a cornerstone technique in analytical chemistry, metabolomics, pharmaceutical research, and many other fields. The prediction of mass spectra from molecular structures, and the inverse inference of molecular structures from spectra, remain active areas of research.
While machine learning methods increasingly improve forward spectrum prediction from known structures, the inverse problem --- predicting plausible molecular structures from a given spectrum --- remains difficult, ambiguous, and poorly interpretable.

This work investigates the integration of auxiliary intermediate fragments generated by graph-transformation tools into machine learning pipelines for bidirectional translation between molecular structures and mass spectra.
These graph-based intermediate representations provide an explicit fragmentation layer between molecules and spectra, making the prediction process more interpretable than purely black-box approaches.

The proposed framework systematically incorporates graph-based fragmentation intermediates into structure $\leftrightarrow$ spectrum translation models.
By combining molecular representations, spectral information, and fragment-level graph structures, the work aims to improve both interpretability and predictive performance in machine learning-based mass spectrometry modeling.

\phantomsection
\addcontentsline{toc}{chapter}{Abstract}
\clearpage
\chapter*{Zusammenfassung}
%min. 500 Zeichen inkl. Leerzeichen
Massenspektrometrie ist eine Schlüsseltechnik in der analytischen Chemie, der Metabolomik, der pharmazeutischen Forschung und vielen weiteren Bereichen.
Die Vorhersage von Massenspektren aus Molekülstrukturen sowie die inverse Ableitung von Molekülstrukturen aus Spektren bleiben aktive Forschungsbereiche.
Während Methoden des maschinellen Lernens die Vorhersage von Spektren aus bekannten Strukturen zunehmend verbessern, bleibt das inverse Problem --- die Vorhersage plausibler Molekülstrukturen aus einem gegebenen Spektrum --- schwierig, mehrdeutig und nur eingeschränkt interpretierbar.

Diese Arbeit untersucht die Integration zusätzlicher Zwischenfragmente, die durch graphbasierte Transformationstools erzeugt werden, in Machine-Learning-Pipelines für die bidirektionale Übersetzung zwischen Molekülstrukturen und Massenspektren.
Diese graphbasierten Zwischenrepräsentationen bilden eine explizite Fragmentierungsschicht zwischen Molekülen und Spektren und machen den Vorhersageprozess dadurch besser interpretierbar als rein Black-Box-basierte Ansätze.

Das vorgeschlagene Framework integriert graphbasierte Fragmentierungszwischenstufen systematisch in Modelle zur Struktur $\leftrightarrow$ Spektrum-Übersetzung.
Durch die Kombination molekularer Repräsentationen, spektraler Information und fragmentbasierter Graphstrukturen zielt die Arbeit darauf ab, sowohl die Interpretierbarkeit als auch die Vorhersageleistung von Machine-Learning-basierten Modellen für die Massenspektrometrie zu verbessern.

\phantomsection
\addcontentsline{toc}{chapter}{Zusammenfassung}
\listoftables
\phantomsection
\addcontentsline{toc}{chapter}{List of Tables}
\listoffigures
\phantomsection
\addcontentsline{toc}{chapter}{List of Figures}
